% Complete IEEEtran manuscript for Overleaf
% Based on the supplied IEEE Transactions journal example.
% Replace the commented journal header only after the target journal is confirmed.

\documentclass[journal]{IEEEtran}

\usepackage{graphicx}
\graphicspath{{figures/}}
\DeclareGraphicsExtensions{.pdf,.png,.jpg,.jpeg}
\usepackage[cmex10]{amsmath}
\usepackage{amssymb}
\usepackage{array}
\usepackage{booktabs}
\usepackage{multirow}
\usepackage[caption=false,font=footnotesize]{subfig}
\usepackage{placeins}
\usepackage{cite}
\usepackage{url}
\usepackage[hidelinks]{hyperref}

\hyphenation{op-tical net-works semi-conduc-tor meta-rein-force-ment}

\newcommand{\clip}{\operatorname{clip}}
\newcommand{\E}{\mathbb{E}}
\newcommand{\eqnote}[1]{\par\vspace{-0.3ex}\noindent #1\par\vspace{0.4ex}}

\begin{document}

\title{A Group Consensus Method for Interval-Valued Fuzzy Preferences Based on Context Encoding and Social Attention}

\author{Yang~Song, Mingshuo~Cao, Zhaoguang~Zhu, Jian~Wu, and Francisco~Chiclana%
\thanks{This work was supported by the National Natural Science Foundation of China under Grant 72601254, the Shanghai Pujiang Talent Program under Grant 24PJC042, and the Shanghai Philosophy and Social Sciences Planning Project under Grant 2025EZH002.}%
\thanks{Yang Song is with the School of Information Engineering, Shanghai Maritime University, Shanghai 201306, China.}%
\thanks{Mingshuo Cao, Zhaoguang Zhu and Jian Wu are with the School of Economics and Management, Shanghai Maritime University, Shanghai 201306, China.}%
\thanks{Francisco Chiclana is with the Institute of Artificial Intelligence, Faculty of Computing, Engineering and Media, De Montfort University, Leicester, U.K.}%
\thanks{Corresponding authors: Mingshuo Cao and Zhaoguang Zhu (e-mail: caomingshuo0606@163.com, zhaoguang-zhu@163.com).}}

% Uncomment and revise only after the target journal is confirmed.
% \markboth{IEEE TRANSACTIONS ON ..., VOL.~XX, NO.~X, 2026}%
% {Song \MakeLowercase{\textit{et al.}}: Interval-Valued Fuzzy Group Consensus Method}

\maketitle

\begin{abstract}
Interval-valued preferences preserve members' acceptable ranges, but they also make consensus control sensitive to endpoint dispersion, hidden response behavior, and changes in group size. This paper develops Interval Meta-GDM as a centralized coordinator that observes the global interval state and interaction history, then issues member-specific suggestions to individuals who respond passively according to unobserved behavioral rules. The coordinator combines equal endpoint translation, a ten-round Transformer context encoder, a variational response code, permutation-invariant social attention, and a soft actor--critic policy. Its reward jointly evaluates endpoint consensus, nonlinear adjustment cost, cost fairness, boundary feasibility, and outlier distance. The model is trained for 5,000 episodes only at $N=10$ and is subsequently frozen. Evaluation uses 50 paired scenarios per condition at $N=10$, 40, and 100 and includes deterministic rules, PPO, DDPG, conventional SAC, interval-width tests, hidden-response tests, context ablation, and interval-preservation checks. In static homogeneous settings, the proposed model and deterministic rules achieve the same final success, so no universal rule advantage is claimed. Under hidden heterogeneity, however, success reaches 0.98, 1.00, and 1.00, compared with 0.80, 0.64, and 0.36 for the strongest rule, while termination occurs 21.92, 45.28, and 72.84 rounds earlier. At $N=40$, all three reinforcement-learning baselines fail within 120 rounds, whereas the proposed model succeeds in 16.72 rounds. Masking history reduces success by 14, 22, and 44 percentage points, and the endpoint-translation test confirms zero mean half-width change. These results support history-conditioned coordination and cross-scale reuse within the tested distributions, while the single training seed and limited scenario count constrain the strength of the conclusions.
\end{abstract}

\begin{IEEEkeywords}
Group decision-making, interval-valued fuzzy preference, consensus, adjustment cost, meta-reinforcement learning.
\end{IEEEkeywords}

\IEEEpeerreviewmaketitle

\section{Introduction}
\label{sec:introduction}

Group consensus models reconcile heterogeneous opinions without treating uncertainty as noise. Fuzzy-set theory supplies a general representation of imprecise assessments \cite{ref1}. Personalized feedback mechanisms and reviews of computing-with-words consensus explain how a moderator can convert disagreement into revision advice \cite{ref6}, \cite{ref25}. Recent TFS studies extend this perspective through trust-consensus multiplex networks, social-trust completeness, and decayed propagation across overlapping communities \cite{tfsWu2022}, \cite{tfsLiu2023}, \cite{tfsJi2024}. These studies are especially relevant because they show that the structure and reliability of social relations affect both the information available to a moderator and the speed at which agreement can be reached. Multiplex trust models distinguish social influence from evolving consensus relations; completeness-aware models identify when missing trust links make feedback unreliable; and decayed propagation limits the influence of long or overlapping trust paths. Their common lesson is that a consensus mechanism should not treat every observed relation as equally informative. Nevertheless, these approaches primarily refine trust propagation or feedback under an explicitly represented network. They do not infer an unobserved threshold, inertia pattern, or acceptance tendency from the discrepancy between a recommendation and the movement that follows. They also do not address the possibility that a pointwise consensus score improves because acceptable intervals have been narrowed. The present study therefore treats network-aware consensus as complementary background and focuses on historical response inference together with endpoint-preserving coordination. Interval-valued preference formulations retain both a central assessment and an admissible uncertainty range \cite{ref16}, \cite{ref17}. Classical opinion dynamics, adaptive feedback, and dynamic-network consensus explain convergence and negotiated adjustment \cite{ref2}, \cite{ref21}, \cite{ref27}, but a point-valued state can still hide endpoint violations or create apparent agreement by suppressing uncertainty.

An interval records two kinds of information. Its center gives the current position, whereas its half-width gives the range the member still accepts. Groups with similar mean centers may therefore have different endpoint dispersion, feasible adjustment space, and boundary risk. A center-only distance cannot determine whether all reported ranges are compatible. The model must use both endpoints in its stopping rule and move the interval without silently narrowing it.

Consensus quality also depends on the path used to reach it. Minimum-cost and fairness-aware models show that the same terminal consensus can distribute burdens very differently \cite{ref7}, \cite{ref8}, \cite{ref20}. Trust-sensitive cost and feedback models further show that a mathematically suitable suggestion may not be fully realized by the recipient \cite{ref12}, \cite{ref26}. A small movement may indicate stubbornness, an unmet threshold, or an inertial response. These explanations look similar in one round but imply different later actions, so repeated recommendations, realized movements, and rewards are needed to infer response tendencies. Minimum-adjustment research further considers informed-individual cost, willingness to modify assessments, and Bayesian trust under uncertain evaluations \cite{tfsLiang2022}, \cite{tfsZhang2022}, \cite{tfsDai2024}. The informed-individual formulation shows that the location of reliable information and the time available for negotiation can change the least-cost adjustment path. Willingness-aware feedback separates a mathematically desirable recommendation from the amount a decision maker is prepared to implement, while Bayesian trust modeling represents uncertainty in interpersonal reliability rather than fixing every trust coefficient in advance. Together, these contributions establish that adjustment cost is inseparable from information quality and behavioral acceptance. They also motivate the distinction used here between a coordinator's proposal and the realized movement. The remaining difficulty is that willingness, threshold, and inertia may be unavailable at deployment and may only become visible through repeated interaction. Instead of assigning those parameters to the policy input, Interval Meta-GDM encodes the observable sequence of proposals, movements, and rewards and uses that history to condition later suggestions. This choice makes the behavioral signal auditable without claiming that the latent code is a recovered psychological label. Recent studies further extend this line through Bayesian minimum adjustment, bounded-rationality modeling, and hybrid crowd--AI decision support \cite{ref28}, \cite{ref29}, \cite{ref30}.

Reinforcement learning can optimize continuous coordination paths in group decision making \cite{ref14}, \cite{ref15}, yet flattened group states bind both input and output dimensions to a fixed group size. Retraining at every $N$ does not demonstrate transfer, and mean pooling can hide influential members. Temporal encoding, permutation-invariant aggregation, and attention offer a possible shared representation across scales, while a context-conditioned latent code can summarize task-relevant behavior from interaction history. The unresolved issue lies at the intersection of four requirements that are often studied separately: the state is set-valued, response characteristics are hidden, group size changes between training and evaluation, and success must be judged jointly with cost and burden distribution. A point-valued controller may satisfy the numerical stopping threshold while violating an interval endpoint. A fixed-length neural input may optimize one training size but require a different architecture when members are added. A policy without interaction history may repeatedly issue ineffective suggestions because threshold and inertial responses look identical in a single observation. Finally, a controller can obtain high success by imposing excessive or highly concentrated adjustment costs. Addressing only one of these issues would therefore leave an alternative explanation for any observed gain. This motivates a centralized, set-based policy whose output dimension is shared across members, whose temporal encoder uses only observable history, and whose evaluation reports endpoint feasibility, success, speed, cost, fairness, and reward together.

These limitations raise four questions. Can continuous control reach consensus without narrowing submitted intervals? Does history help when response parameters are hidden from every method? Can one checkpoint trained at $N=10$ operate at $N=40$ and $N=100$ without architectural change? Finally, do any gains remain when success, rounds, nonlinear cost, burden fairness, reward, and boundary violations are reported together? Separating these questions prevents a high reward at one scale from being mistaken for evidence of interval fidelity, historical adaptation, and structural transfer.

The proposed Interval Meta-GDM combines a ten-round Transformer, a variational response code, social attention, and a shared SAC actor. Equal endpoint translation preserves half-width, while the reward uses endpoint consensus, nonlinear cost, fairness, outlier distance, and boundary violations. Evaluation proceeds from static rules and fixed-dimensional PPO, DDPG, and SAC to interval-width, unseen-scale, hidden-heterogeneity, history-masking, and Hausdorff-consistency tests. This order distinguishes basic solvability from the conditions in which context and set structure provide an advantage.

The contributions are threefold. First, centers, half-widths, endpoints, boundaries, and costs are defined consistently, and actions translate rather than shrink intervals. Second, history encoding and social attention convert observed interactions into member-specific actions with shared parameters. Third, frozen-policy experiments compare rule and learning baselines, scale transfer, hidden responses, and context ablation using the same metric set and paired scenarios.

The remainder of this paper is organized as follows. Section~II formulates interval-valued preferences, endpoint consensus, hidden responses, interval translation, and adjustment cost. Section~III presents the centralized Interval Meta-GDM coordinator, including context encoding, social attention, reward construction, SAC optimization, and training behavior. Section~IV describes the paired evaluation protocol and reports rule and reinforcement-learning comparisons, interval and scale tests, hidden-heterogeneity results, context ablation, sensitivity analysis, and process-level reproducibility evidence. Section~V concludes the paper and identifies the principal limitations and directions for further validation.
\section{Problem Formulation}
\label{sec:problem}

The formulation follows one complete negotiation round. Individual intervals first provide a common representation of position and uncertainty \cite{ref16}, \cite{ref17}. Their endpoints are then used to decide whether the group is both concentrated and feasible. A recommendation subsequently passes through an unobserved response, updates the interval without changing its width, and produces cost and fairness signals. This order matters because the policy controls recommendations, but not their acceptance.

\subsection{Interval Preference and Consensus}

The calculation begins at the member level. For $N$ members, $I_{i,t}$ denotes member $i$'s acceptable interval at round $t$; $h_{i,t}=0$ recovers the point-valued case.

\begin{equation}
I_{i,t}=[l_{i,t},u_{i,t}]
=[c_{i,t}-h_{i,t},c_{i,t}+h_{i,t}],\quad
0\le l_{i,t}\le u_{i,t}\le 1.
\label{eq:interval}
\end{equation}

\eqnote{Eq.~\eqref{eq:interval} separates the negotiable center $c_{i,t}$ from the uncertainty half-width $h_{i,t}$ that must be retained. Its two endpoints become the common representation used by all later consensus and boundary calculations.}

\begin{equation}
\begin{aligned}
\bar l_t&=N^{-1}\sum_i l_{i,t}, &
\bar u_t&=N^{-1}\sum_i u_{i,t},\\
\bar c_t&=N^{-1}\sum_i c_{i,t}, &
\bar h_t&=N^{-1}\sum_i h_{i,t}.
\end{aligned}
\label{eq:means}
\end{equation}

\eqnote{Once every interval has the same representation, the averages in Eq.~\eqref{eq:means} establish a group reference. They are not a consensus result themselves; they are the baselines against which endpoint and member deviations are measured next.}

These group references make it possible to measure disagreement, but a valid stopping rule must also reject endpoint violations. Concentration and feasibility are therefore evaluated separately before they are combined \cite{ref6}, \cite{ref21}, \cite{ref27}. The cutoff $\delta=0.01$ distinguishes negligible endpoint dispersion from dispersion that still requires penalization, whereas $\kappa$ controls the exponential decay rate; the default experiments use $\kappa=5$ and $C^*=0.90$.

\begin{equation}
D_t=\left\{\frac{1}{2N}\sum_i\left[
(l_{i,t}-\bar l_t)^2+(u_{i,t}-\bar u_t)^2
\right]\right\}^{1/2}.
\label{eq:dispersion}
\end{equation}

\eqnote{Using those endpoint averages, $D_t$ measures geometric dispersion and vanishes only when all lower and upper endpoints coincide. Dispersion alone, however, cannot reject a tightly clustered but infeasible group, which motivates the safety test.}

\begin{equation}
\begin{aligned}
p_t^-&=\max\left\{0,\frac{s_L-\min_i l_{i,t}}{s_L}\right\},\\
p_t^+&=\max\left\{0,\frac{\max_i u_{i,t}-s_U}{1-s_U}\right\},\\
S_t&=\clip\!\left[1-\max(p_t^-,p_t^+),0,1\right].
\end{aligned}
\label{eq:safety}
\end{equation}

\eqnote{The factors $p_t^-$ and $p_t^+$ measure violations below $s_L$ and above $s_U$. When the absolute lower bound is $s_L=0$, the implementation replaces the denominator of $p_t^-$ by $\max(s_L,\varepsilon_b)$, with $\varepsilon_b=10^{-8}$; equivalently, a conditional branch returns zero when all lower endpoints are nonnegative and uses the stabilized denominator otherwise. The same safeguard is applied to $1-s_U$ when $s_U=1$. Their maximum determines $S_t$, so a cluster outside the admissible region cannot be mistaken for valid consensus.}

\begin{equation}
F(D_t)=
\begin{cases}
1, & D_t<\delta,\\
\exp(-\kappa D_t), & D_t\ge\delta,
\end{cases}
\qquad C_t=F(D_t)S_t.
\label{eq:consensus}
\end{equation}

\eqnote{Eq.~\eqref{eq:consensus} combines the two tests. The nondecaying branch applies only when $D_t<\delta$; for $D_t\ge\delta$, the factor $\exp(-\kappa D_t)$ penalizes remaining dispersion. Thus, $\delta$ selects the branch and $\kappa$ controls the strength of consensus decay, consistently with Table~\ref{tab:sensitivity}. The product $C_t=F(D_t)S_t$ is the stopping signal, and a run succeeds only when $C_t\ge C^*$.}

\begin{equation}
\begin{aligned}
d_{i,t}&=\left[\frac{(l_{i,t}-\bar l_t)^2+(u_{i,t}-\bar u_t)^2}{2}\right]^{1/2},\\
v_{i,t}&=\max(s_L-l_{i,t},0)+\max(u_{i,t}-s_U,0).
\end{aligned}
\label{eq:member_distance}
\end{equation}

\eqnote{The global score determines whether the group may stop, but the policy also needs member-level feedback. Eq.~\eqref{eq:member_distance} supplies endpoint deviation $d_{i,t}$ and boundary excess $v_{i,t}$, which later enter each local reward.}

\subsection{Response, Interval Update, and Cost}

The second part of the round distinguishes advice from behavior \cite{ref18}, \cite{ref23}, \cite{ref24}. The policy proposes $a_{i,t}\in[-1,1]$, but hidden responses determine what is actually accepted. The accepted movement then updates the interval and creates an adjustment burden. Experiments use $\zeta=10$, $\lambda=0.35$, and $\Delta=0.1$.

\begin{equation}
m_{i,t}=\clip\!\left[(1-\omega_i)\Delta
f_i(a_{i,t},a_{i,t-1};\theta_i),-\Delta,\Delta\right].
\label{eq:movement}
\end{equation}

\eqnote{The actor's proposal $a_{i,t}$ does not move an interval directly. Eq.~\eqref{eq:movement} first converts it into realized movement $m_{i,t}$ through willingness $1-\omega_i$, response function $f_i$, and the limit $\Delta$.}

\begin{equation}
\begin{aligned}
f_0&=a_{i,t},\\
f_1&=a_{i,t}\mathbf{1}(|a_{i,t}|\ge\theta_i),\\
f_2&=a_{i,t}\sigma\!\left[\zeta(|a_{i,t}|-\theta_i)\right],\\
f_3&=\lambda a_{i,t}+(1-\lambda)a_{i,t-1}.
\end{aligned}
\label{eq:responses}
\end{equation}

\eqnote{The four definitions of $f_i$ produce linear, hard-threshold, smooth-threshold, and inertial reactions. Because their type and parameters are hidden, the gap between proposed and realized movement becomes evidence for the history encoder.}

\begin{equation}
c_{i,t+1}=\clip[c_{i,t}+m_{i,t},h_{i,t},1-h_{i,t}],
\qquad h_{i,t+1}=h_{i,t}.
\label{eq:interval_update}
\end{equation}

\eqnote{After the hidden response is applied, Eq.~\eqref{eq:interval_update} updates only the center and clips it to a feasible range. Keeping $h_{i,t+1}=h_{i,t}$ translates both endpoints equally, so consensus is never obtained by erasing uncertainty.}

Quadratic, fairness-aware, and dual-behavior cost models motivate this design \cite{ref9}, \cite{ref10}, \cite{ref11}. Utility-based preference learning supports the burden distinction \cite{ref13}. Reported costs use $\eta=75$ and $p=2$; the shifted Gini measures how evenly those costs are distributed.

\begin{equation}
J_{i,t}=q_i\eta\Delta\left(\frac{|m_{i,t}|}{\Delta}\right)^p,
\qquad p>1.
\label{eq:cost}
\end{equation}

\eqnote{Realized movement also creates a burden. Eq.~\eqref{eq:cost} weights it by sensitivity $q_i$, scales it through $\eta$, and uses $p>1$ so a large adjustment costs more than several smaller ones.}

\begin{equation}
\begin{aligned}
\widetilde J_{i,t}&=J_{i,t}-\min_j J_{j,t},\\
g_t&=
\begin{cases}
0, & \sum_i\widetilde J_{i,t}=0,\\[2pt]
\displaystyle
\frac{2\sum_i i\widetilde J_{(i),t}}
{N\sum_i\widetilde J_{i,t}}-\frac{N+1}{N}, & \text{otherwise}.
\end{cases}
\end{aligned}
\label{eq:gini}
\end{equation}

\eqnote{Total cost does not reveal who bears it. The shifted and ordered costs in Eq.~\eqref{eq:gini} produce $g_t$, which is zero under equal burdens and increases with concentration; this fairness signal completes the environment description.}

\section{Interval Meta-GDM}
\label{sec:model}

\subsection{Architecture and Context Encoding}

Interval Meta-GDM implements the preceding round as a centralized coordination loop rather than a multi-agent reinforcement-learning system. At round $t$, one coordinator observes the complete group state as a set of interval-valued member states and receives the observable histories generated by earlier interactions. Its context encoder infers member-specific response tendencies, social attention summarizes the group configuration, and one shared central policy produces the suggestion vector $a_t=(a_{1,t},\ldots,a_{N,t})$. Suggestions may differ across members, but they are generated jointly by the same coordinator.

The individuals are passive response units rather than learning agents. They do not maintain policies, optimize rewards, exchange learned messages, or update network parameters. Instead, each member reacts to suggestions issued by the coordinator according to hidden response characteristics, after which the environment returns consensus, cost, fairness, and boundary feedback to the centralized learner. During training, these transitions update the encoder, actor, and twin critics of the coordinator through replay; testing freezes the entire mapping. Fig.~\ref{fig:architecture} separates the interaction and centralized update flows.

\begin{figure*}[!t]
\centering
\includegraphics[width=0.98\textwidth]{fig1_architecture_mentor.png}
\caption{Interaction decision and SAC training-update loops of Interval Meta-GDM.}
\label{fig:architecture}
\end{figure*}

The interaction loop begins by separating what is directly observable from what must be inferred. A five-dimensional state describes the current interval and group position, whereas $L=10$ interaction tokens describe how the member has reacted over time. The Transformer \cite{ref3} turns the latter evidence into a response code.

\begin{equation}
s_{i,t}=\left[c_{i,t},h_{i,t},c_{i,t}-\bar c_t,
h_{i,t}-\bar h_t,\frac{H-t}{H}\right]\in\mathbb{R}^5.
\label{eq:state}
\end{equation}

\eqnote{The state in Eq.~\eqref{eq:state} records absolute center and width, their deviations from the group, and remaining time. It explains where a member is now, but not how that member has reacted to earlier advice.}

\begin{equation}
\begin{aligned}
x_{i,t}&=\left[a_{i,t},m_{i,t},\frac{r_{i,t}}{\nu}\right],\\
\mathcal H_{i,t}&=(x_{i,t-L},\ldots,x_{i,t-1}),\\
\mathcal H_{i,t}&\in\mathbb{R}^{L\times3}.
\end{aligned}
\label{eq:history}
\end{equation}

\eqnote{That missing behavioral evidence is supplied by Eq.~\eqref{eq:history}. Each token aligns proposal, realized movement, and reward, and the latest $L$ tokens form $\mathcal H_{i,t}$, exposing persistent thresholds or inertia.}

\begin{equation}
q_\phi(z_{i,t}\mid\mathcal H_{i,t})=
\mathcal N\!\left(\mu_\phi(\mathcal H_{i,t}),
\operatorname{diag}[\sigma_\phi^2(\mathcal H_{i,t})]\right).
\label{eq:posterior}
\end{equation}

\eqnote{The Transformer maps this history to the posterior in Eq.~\eqref{eq:posterior}. Its mean and diagonal variance represent uncertainty over latent response characteristics instead of assigning an observed personality label \cite{ref31}, \cite{ref34}.}

\begin{equation}
\begin{aligned}
z_{i,t}&=\mu_{i,t}+\sigma_{i,t}\odot\epsilon,
&\epsilon&\sim\mathcal N(0,I),\\
\chi_{i,t}&=\frac12\sum_k\left(\mu_{i,k}^2+\sigma_{i,k}^2
-\log\sigma_{i,k}^2-1\right),
&z_{i,t}^{\mathrm{test}}&=\mu_{i,t}.
\end{aligned}
\label{eq:latent}
\end{equation}

\eqnote{Eq.~\eqref{eq:latent} makes the posterior trainable through reparameterization and regularizes it with $\chi_{i,t}$. Training samples $z_{i,t}$, whereas testing uses its mean so the frozen checkpoint acts deterministically.}

\subsection{Social Attention and Shared Policy}

Within the centralized coordinator, state-code features are fused and pooled by four-head social attention to determine how each member relates to the rest of the group. The resulting group context conditions the shared squashed-Gaussian actor, while the central twin critics evaluate the joint consequence. Social attention is therefore an internal aggregation mechanism of the coordinator, not a communication protocol among autonomous agents.

\begin{equation}
u_{i,t}=\operatorname{ReLU}(W_x[s_{i,t};z_{i,t}]+b_x),
\qquad \bar u_t=N^{-1}\sum_i u_{i,t}.
\label{eq:feature}
\end{equation}

\eqnote{The observable state and inferred response code meet in Eq.~\eqref{eq:feature}. Feature $u_{i,t}$ describes one member, while the mean $\bar u_t$ provides a group anchor whose dimension is independent of $N$.}

\begin{equation}
\begin{aligned}
q_{h,t}&=W_{q,h}\bar u_t+b_{q,h},\\
k_{i,h,t}&=W_{k,h}u_{i,t},\qquad
v_{i,h,t}=W_{v,h}u_{i,t},\\
\alpha_{i,h,t}&=\operatorname{softmax}_i
\left(\frac{q_{h,t}^{\mathsf T}k_{i,h,t}}{\sqrt{d_h}}\right).
\end{aligned}
\label{eq:attention}
\end{equation}

\eqnote{Starting from that anchor, Eq.~\eqref{eq:attention} compares a group query with every member key. The normalized weights $\alpha_{i,h,t}$ identify the members most relevant to the current decision instead of treating all deviations equally.}

\begin{equation}
o_t=W_o\operatorname{LN}\!\left(
\operatorname*{Concat}_h\left[\sum_i\alpha_{i,h,t}v_{i,h,t}+\bar u_t\right]
\right).
\label{eq:attention_output}
\end{equation}

\eqnote{Eq.~\eqref{eq:attention_output} turns those weights into group context $o_t$. Weighted values, the residual group mean, normalization, and head concatenation preserve salient members while remaining invariant to input order \cite{ref4}, \cite{ref19}.}

\begin{equation}
\begin{aligned}
(\widehat\mu_{i,t},\log\widehat\sigma_{i,t})
&=f_\theta(s_{i,t},z_{i,t},o_t),\\
\xi_{i,t}&=\widehat\mu_{i,t}+\widehat\sigma_{i,t}\epsilon_i,\\
a_{i,t}&=\tanh(\xi_{i,t}).
\end{aligned}
\label{eq:actor}
\end{equation}

\eqnote{The shared actor in Eq.~\eqref{eq:actor} combines member state, latent response code, and group context. It samples a Gaussian pre-action for exploration and applies $\tanh$ so each suggestion remains in $[-1,1]$.}

\begin{equation}
\begin{aligned}
&\log\pi_\theta(a_{i,t}\mid s_{i,t},z_{i,t},o_t)\\
&\quad=\log\mathcal N(\xi_{i,t};\widehat\mu_{i,t},\widehat\sigma_{i,t}^2)\\
&\qquad-\log(1-a_{i,t}^2+\epsilon_0).
\end{aligned}
\label{eq:logprob}
\end{equation}

\eqnote{Because the action is squashed, its probability cannot use the Gaussian density alone. Eq.~\eqref{eq:logprob} adds the $\tanh$ Jacobian correction and supplies the valid entropy term used by the SAC objectives.}

\subsection{Reward and Optimization}

The action becomes useful only when its consequence is translated back into the objectives used to define consensus. The reward therefore links group progress to individual burden and boundary safety before the transition enters replay. Reward weights are $w_C=50$, $b=10$, $w_F=1$, $\tau=0.1$, $w_D=2$, and $w_B=20$, with $R=20$.

\begin{equation}
\begin{aligned}
r_{i,t}=\clip\big[&w_C(C_{t+1}-C_t)
+b\mathbf{1}(C_{t+1}\ge C^*)-w_Fg_t-\tau\\
&-w_Dd_{i,t}-w_Bv_{i,t}-J_{i,t},-R,R\big].
\end{aligned}
\label{eq:reward}
\end{equation}

\eqnote{The reward in Eq.~\eqref{eq:reward} closes the interaction loop. Consensus progress and terminal success are positive; delay, unfairness, outlier distance, boundary excess, and cost are negative. Clipping prevents isolated rounds from destabilizing value learning.}

The reward coefficients were selected through empirical tuning in the $N=10$ training environment and were fixed before the frozen-policy evaluations. Their different magnitudes compensate for differences in the natural scales of the components rather than representing universal preference weights. The per-round consensus increment is usually small and is therefore amplified by $w_C=50$, whereas the safety-critical boundary term receives $w_B=20$ and the terminal bonus $b=10$ distinguishes successful completion. The fairness, outlier-distance, time, and adjustment-cost terms retain smaller or naturally scaled contributions. These values should therefore be interpreted as empirically calibrated settings, not theoretically optimal constants.

The history length $L=10$ follows the same practical criterion. A shorter window may omit repeated proposal--movement patterns needed to distinguish threshold and inertial responses, whereas a substantially longer sequence increases Transformer computation, replay-memory use, and exposure to less relevant early-round information. Ten rounds provide sufficient recent behavioral evidence while keeping centralized inference and off-policy training manageable.

\begin{equation}
\bar r_t=N^{-1}\sum_i r_{i,t},\qquad
(s_t,\mathcal H_t,a_t,\bar r_t,s_{t+1},\mathcal H_{t+1},d_t)\rightarrow\mathcal D.
\label{eq:replay}
\end{equation}

\eqnote{Eq.~\eqref{eq:replay} then turns one interaction round into an off-policy sample. Both current and next histories are stored because the latent code must be reconstructed consistently on both sides of the Bellman update.}

Training uses $\gamma=0.95$, $\alpha=0.1$, $\tau_q=0.01$, $\beta_z=0.05$, learning rate $3\times10^{-4}$, and at most $H=120$ rounds before replay updates.

\begin{equation}
\begin{aligned}
y_t={}&\bar r_t+\gamma(1-d_t)\Bigl[
\min_{j\in\{1,2\}}\bar Q_j(s_{t+1},z',a')\\
&\hspace{18mm}-\alpha\log\pi_\theta(a'\mid s_{t+1},z')\Bigr].
\end{aligned}
\label{eq:sac_target}
\end{equation}

\eqnote{Following soft actor--critic \cite{ref5}, Eq.~\eqref{eq:sac_target} builds the target from immediate reward, the smaller target-critic estimate, and policy entropy. The factor $1-d_t$ stops bootstrapping at a terminal state.}

\begin{equation}
\mathcal L_Q=\E_{\mathcal B}\left[
\sum_{j=1}^{2}\left(Q_j(s_t,z_t,a_t)-y_t\right)^2\right].
\label{eq:critic_loss}
\end{equation}

\eqnote{The critic loss in Eq.~\eqref{eq:critic_loss} fits both value estimates to that target. Two critics are retained because their conservative minimum is used in target construction and policy improvement to limit overestimation.}

\begin{equation}
\mathcal L_\pi=\E_{\mathcal B}\left[
\alpha\log\pi_\theta(a_t\mid s_t,z_t)
-\min_{j\in\{1,2\}}Q_j(s_t,z_t,a_t)\right].
\label{eq:actor_loss}
\end{equation}

\eqnote{With the critics updated, Eq.~\eqref{eq:actor_loss} favors actions with high conservative value while retaining entropy. It may legitimately become negative when the value term dominates, so its sign is not a convergence criterion.}

\begin{equation}
\mathcal L_z=\mathcal L_Q+\beta_zN^{-1}\sum_i
\operatorname{KL}\!\left[q_\phi(z_{i,t}\mid\mathcal H_{i,t})
\,\|\,\mathcal N(0,I)\right].
\label{eq:encoder_loss}
\end{equation}

\eqnote{The encoder is trained through Eq.~\eqref{eq:encoder_loss}. Critic error makes the latent code useful for decisions, while the $\beta_z$-weighted KL term prevents unconstrained memorization; no response labels are required.}

\begin{equation}
\bar\psi_j\leftarrow(1-\tau_q)\bar\psi_j+\tau_q\psi_j,
\qquad j\in\{1,2\}.
\label{eq:soft_update}
\end{equation}

\eqnote{Finally, Eq.~\eqref{eq:soft_update} moves each target critic slowly toward its online counterpart. This Polyak update completes one SAC cycle and prevents the next Bellman targets from changing as abruptly as the online networks.}

The full chain is now closed: interval observations and realized responses form the context; context and social attention produce actions; actions change the environment; and the resulting reward updates the same representation on the next replay step. The training curves should therefore be interpreted as coupled evidence about this loop, not as four independent objectives.

\subsection{Training Behavior}

Fig.~\ref{fig:training} reports the audited 5,000-episode history. Reward and terminal consensus improve after early exploration and stabilize despite curriculum changes. Actor loss becomes negative as high-value actions dominate the entropy term; it is not expected to decrease monotonically. Critic loss is initially large because targets change with the policy, then declines and fluctuates around a bounded level. Taken together, the curves show policy improvement rather than a divergent objective.

\begin{figure}[!t]
\centering
\subfloat[Episode reward]{\includegraphics[width=0.94\columnwidth]{fig02a_episode_reward.png}}\par\vspace{-0.6ex}
\subfloat[Terminal consensus]{\includegraphics[width=0.94\columnwidth]{fig02b_final_consensus.png}}\par\vspace{-0.6ex}
\subfloat[Actor loss]{\includegraphics[width=0.94\columnwidth]{fig02c_actor_loss.png}}\par\vspace{-0.6ex}
\subfloat[Critic loss]{\includegraphics[width=0.94\columnwidth]{fig02d_critic_loss.png}}
\caption{Training diagnostics: (a) reward; (b) terminal consensus; (c) actor loss; and (d) critic loss.}
\label{fig:training}
\end{figure}

\section{Experiments}
\label{sec:experiments}

\subsection{Protocol and Metrics}

The main model uses seed 20260507 and is trained only at $N=10$ for 5,000 episodes. Frozen evaluation uses 50 paired scenarios per cell at $N=10$, 40, and 100, with identical initial intervals and hidden parameters across methods. PPO, DDPG, and conventional SAC are trained separately for each tested $N$. No replay, backpropagation, or target update occurs during testing.

The principal experimental, architectural, and optimization settings are summarized in Table~\ref{tab:settings}.

In the hidden-heterogeneity experiments, each member is independently assigned one of the four response types in Eq.~\eqref{eq:responses}: linear $f_0$, hard-threshold $f_1$, smooth-threshold $f_2$, or inertial $f_3$. The assignment follows a discrete uniform distribution, so every type has probability $0.25$ at $N=10$, 40, and 100. Realized counts are not forced to be exactly equal, particularly at $N=10$, but approach an even composition in expectation. Response types are resampled for each scenario using the recorded scenario seed and remain fixed throughout the episode; threshold parameters are sampled as $\theta_i\sim\mathcal U(0.18,0.55)$. All methods in a paired scenario receive the same response assignments and hidden parameters, but neither the central policy nor the rule baselines can observe them directly. The static control setting uses only the linear response $f_0$.

\begin{table*}[!t]
\caption{Key Experimental and Training Settings}
\label{tab:settings}
\centering
\footnotesize
\setlength{\tabcolsep}{5pt}
\begin{tabular}{llll}
\toprule
Module & Setting & Symbol & Value \\
\midrule
Environment & Training/test scale & $N$ & 10 / 10, 40, 100 \\
Environment & Episodes/scenarios & -- & 5,000 / 50 per cell \\
Environment & Horizon/threshold & $H/C^*$ & 120 / 0.90 \\
Intervals & Half-width/cutoff & $h/\delta$ & $[0.02,0.08]$ / 0.01 \\
Context & Window/latent size & $L/d_z$ & 10 / 8 \\
Network & Attention & $d_g$/heads & 256 / 4 \\
SAC & Learning rate/discount & $\mathrm{lr}/\gamma$ & $3\times10^{-4}$ / 0.95 \\
SAC & Soft update/entropy & $\tau_q/\alpha$ & 0.01 / 0.10 \\
Cost & Scale/exponent & $\eta/p$ & 75 / 2 \\
\bottomrule
\end{tabular}
\end{table*}

All metrics use paired frozen-policy episodes. The 95\% intervals use $B=10{,}000$ paired bootstrap resamples; cost is interpreted jointly with success.

\begin{equation}
\begin{aligned}
p_s&=M^{-1}\sum_e\mathbf{1}_e,\\
\bar T&=M^{-1}\sum_eT_e,\\
\bar J&=M^{-1}\sum_e\sum_t\sum_iJ_{i,e,t},\\
\bar R&=M^{-1}\sum_e\sum_t\left(N^{-1}\sum_ir_{i,e,t}\right),\\
\bar V&=M^{-1}\sum_e\sum_t\sum_i
\mathbf{1}(l_{i,e,t}<s_L\ \text{or}\ u_{i,e,t}>s_U).
\end{aligned}
\label{eq:metrics}
\end{equation}

\eqnote{Over $M$ episodes, these terms report success, rounds, total cost, per-member mean reward, and boundary counts. Their joint use prevents failed inactive policies from appearing favorable through low cost alone.}

\begin{equation}
A_{s,e}=\mathbf{1}_e\frac{H-T_e+1}{H+1},\qquad
A_{c,e}=H^{-1}\sum_{t=1}^{H}\widetilde C_{e,t}.
\label{eq:auc}
\end{equation}

\eqnote{The score $A_{s,e}$ is zero for failure and rewards earlier success. The consensus AUC averages the padded trajectory, measuring progress even when methods share the same final success label.}

\begin{equation}
\begin{aligned}
\delta_e&=m_{1,e}-m_{0,e},\qquad
\bar\delta=M^{-1}\sum_e\delta_e,\\
q_L&=Q_{\vartheta/2}(\{\bar\delta_b^*\}_{b=1}^{B}),\\
q_U&=Q_{1-\vartheta/2}(\{\bar\delta_b^*\}_{b=1}^{B}),\\
\mathcal I_{1-\vartheta}&=[q_L,q_U].
\end{aligned}
\label{eq:bootstrap}
\end{equation}

\eqnote{Each $\delta_e$ is a within-scenario model difference. Resampling paired differences preserves scenario matching, and quantiles $q_L,q_U$ form the $(1-\vartheta)$ interval used to assess stability.}

Code, checkpoints, scenario-level data, and all data-driven figure scripts are available at \url{https://github.com/songyang531/IntervalFuzzy-MetaGDM}; see \texttt{scripts/generate\_paper\_figures.py} and \texttt{docs/PAPER\_EXPERIMENT\_MAP.md}.

\subsection{Sensitivity and Reproducibility Settings}

The sensitivity study reuses the frozen evaluation protocol rather than selecting settings from test outcomes. The checkpoint is trained for 5,000 episodes at $N=10$ with seed 20260507, $L=10$, $d_z=8$, and four attention heads. All networks and normalization constants are frozen for evaluation. Each cell contains 50 paired scenarios with identical initial intervals and hidden responses. Table~\ref{tab:sensitivity} reports the tested ranges and the observed boundaries that matter for interpreting the main experiments.

\begin{table*}[!t]
\caption{Training Settings, Sensitivity Ranges, and Observed Boundaries}
\label{tab:sensitivity}
\centering
\scriptsize
\setlength{\tabcolsep}{3pt}
\begin{tabular}{p{0.14\textwidth}p{0.20\textwidth}p{0.58\textwidth}}
\toprule
Item & Tested setting & Observed evidence and interpretation \\
\midrule
Optimization & Learning rate $3\times10^{-4}$; $\gamma=0.95$; hidden dimension 256; $\tau_q=0.01$ & PPO uses clip 0.2. DDPG and SAC use batch 64 and replay 100000; Interval Meta-GDM uses entropy 0.1, batch 32, and replay 100000. \\
Latent code $z$ & $\{4,8,16,32,64\}$; default 8 & At $N=100$, $z=8$ needs 19.02 rounds with reward 18.86; $z=32$ gives 28.80/11.09 and $z=64$ gives 31.14/8.99. Larger codes remain successful but are more conservative. \\
Cost coefficient $\eta$ & $\{25,50,75,100,125\}$; default 75 & At $N=100$, mean rounds are 36.86, 18.66, 19.02, 19.36, and 21.06. The middle region is stable, but the relation is not monotonic. \\
Consensus gain $w_C$ & $\{30,40,50,60,70\}$; default 50 & At $N=100$, $w_C=50$ gives 19.08 rounds, cost 284.64, and reward 19.49; $w_C=70$ gives 33.94, 647.61, and 14.18. \\
Consensus decay $\kappa$ & $\{3,5,7\}$; default 5 & Success at $N=10/40/100$ is 1.00/1.00/1.00, 0.98/1.00/1.00, and 0.36/0.06/0.00. At $\kappa=7$, the $N=100$ center SD is 0.0006 but interval consensus is 0.8857. \\
\bottomrule
\end{tabular}
\end{table*}

The sensitivity pattern is interpretable rather than monotonic. Larger latent codes slow coordination, $\eta=75$ lies in a stable middle region, and excessive consensus gain raises movement cost. When $\kappa=7$, centers nearly coincide but preserved half-width differences keep endpoint consensus below the target. This is an interval-reachability boundary, not evidence that the optimizer diverged.

\subsection{Static Rules and Reinforcement-Learning Baselines}

The deterministic controls are mean convergence, bounded confidence, and observable cost awareness, defined by Eq.~\eqref{eq:rule_mean} through Eq.~\eqref{eq:rule_cost}; a random policy is included as a failure control. In the static identically distributed setting, all deterministic rules and Interval Meta-GDM reach 1.00 success at every $N$. Rules are faster and cheaper at $N=10$. At $N=40$ and 100, however, the model obtains mean rewards of 20.41 and 19.65 and cost Gini values of 0.5292 and 0.5326, the best composite reward and burden balance among deterministic methods. Thus, static tests establish competence, not a universal advantage over rules.

\begin{equation}
a_{i,t}^{(1)}=\clip\!\left[\frac{\bar c_t-c_{i,t}}{\Delta},-1,1\right].
\label{eq:rule_mean}
\end{equation}

\eqnote{This rule moves directly toward the current mean center. Division by $\Delta$ converts displacement to the learned policy's action scale, and clipping applies the common bound.}

\begin{equation}
\begin{aligned}
\mathcal N_{i,t}&=\{j:|c_{j,t}-c_{i,t}|\le\varepsilon\},\\
c_{i,t}^{(2)*}&=\frac{\beta}{|\mathcal N_{i,t}|}
\sum_{j\in\mathcal N_{i,t}}c_{j,t}+(1-\beta)\bar c_t,\\
a_{i,t}^{(2)}&=\clip\!\left[\frac{c_{i,t}^{(2)*}-c_{i,t}}{\Delta},-1,1\right].
\end{aligned}
\label{eq:rule_bounded}
\end{equation}

\eqnote{The set $\mathcal N_{i,t}$ contains neighbors within radius $\varepsilon$. Parameter $\beta$ blends their local mean with the global center before the target gap becomes a clipped action.}

\begin{equation}
\begin{aligned}
\widehat b_{i,t}&=\operatorname{mean}_k\max(-r_{i,k},0),
&\widehat\rho_{i,t}&=\operatorname{mean}_k\frac{|m_{i,k}|}{\Delta},\\
\widehat q_{i,t}&=1+\widehat b_{i,t}/s,
&w_{i,t}&=1/\widehat q_{i,t},\\
c_{w,t}&=\frac{\sum_iw_{i,t}c_{i,t}}{\sum_iw_{i,t}},\\
\lambda_{i,t}&=\clip\!\left[
\frac{\lambda_0+\widehat\rho_{i,t}}{\widehat q_{i,t}},\lambda_m,1\right],\\
a_{i,t}^{(3)}&=\clip\!\left[
\frac{\lambda_{i,t}(c_{w,t}-c_{i,t})}{\Delta},-1,1\right].
\end{aligned}
\label{eq:rule_cost}
\end{equation}

\eqnote{Past negative burden and response rate produce proxy sensitivity $\widehat q_{i,t}$ and inverse-cost center $c_{w,t}$. The adaptive $\lambda_{i,t}$ reduces movement for costly or unresponsive members.}

\begin{figure}[!b]
\centering
\subfloat[Success]{\includegraphics[width=0.72\columnwidth]{fig03a_success_rate.png}}\par\vspace{-1.2ex}
\subfloat[Rounds]{\includegraphics[width=0.72\columnwidth]{fig03b_convergence_rounds.png}}\par\vspace{-1.2ex}
\subfloat[Cost]{\includegraphics[width=0.72\columnwidth]{fig03c_adjustment_cost.png}}\par\vspace{-1.2ex}
\subfloat[Reward]{\includegraphics[width=0.72\columnwidth]{fig03d_episode_reward.png}}
\caption{Reinforcement-learning baselines: (a) success; (b) rounds; (c) cost; and (d) reward.}
\label{fig:rl_baselines}
\end{figure}

Fig.~\ref{fig:rl_baselines} compares fixed-dimensional PPO, DDPG, and SAC baselines. At $N=10$, the proposed model reaches 0.96 success in 15.12 rounds, versus 0.94/20.40 for PPO and 0.82/32.60 and 0.82/34.04 for DDPG and SAC. At $N=40$, all three baselines fail within 120 rounds, whereas Interval Meta-GDM reaches 1.00 success in 16.72 rounds with no boundary violations. Low cost from a failed baseline is not treated as an advantage. This comparison is intended to isolate whether a learned continuous controller can remain operational when the evaluation size changes; it is not evidence that PPO, DDPG, or SAC is intrinsically unsuitable for group decision making. The conventional baselines retain fixed-dimensional state-action mappings, whereas the proposed coordinator applies the same member encoder, attention aggregation, and shared actor at every $N$. The $N=40$ result therefore supports the architectural claim of size-compatible reuse under the stated training and testing protocol. It also explains why success, rounds, and violations are reported before cost: a controller that stops moving and never reaches consensus can appear inexpensive only because it has not solved the task. Conversely, the lower round count of Interval Meta-GDM matters because the interaction horizon is itself a negotiation burden, even when its monetary cost is not represented explicitly.

\subsection{Interval Fuzziness, Scale Transfer, and Hidden Heterogeneity}

Point-valued $[0,0]$ and narrow-interval $[0.01,0.04]$ tests reach 1.00 success at $N=10$, 40, and 100, showing compatibility with precise and mildly uncertain preferences. Under the training interval $[0.02,0.08]$, success is 0.96, 1.00, and 1.00; AUC is 0.8698, 0.8392, and 0.8417. The checkpoint trained at $N=10$ is reused unchanged at $N=40$ and 100, where it requires 19.46 and 19.16 rounds. Here scale denotes group size rather than information granularity \cite{ref22}. The small efficiency loss without retraining supports cross-scale reuse within the tested distributions. This result follows from the way scale enters the coordinator rather than from padding or truncating a fixed vector. Each member is encoded by the same functions, attention weights are normalized over the members present in the current group, and the actor reuses one parameterization when producing member-specific suggestions. Increasing $N$ therefore changes the number of encoded elements but not the dimensionality of an individual feature or the group representation. The result should nevertheless be read as interpolation over the tested task family, not as unrestricted size invariance: the initialization ranges, response mechanisms, safety bounds, and horizon remain those defined by the evaluation protocol. Testing much larger groups, different network structures, or wider intervals would require a new validation rather than extrapolation from the three reported values.

Hidden-heterogeneity tests vary stubbornness, cost sensitivity, and response type while withholding those parameters from every method. Fig.~\ref{fig:heterogeneity} shows that Interval Meta-GDM reaches 0.98, 1.00, and 1.00 success, compared with 0.80, 0.64, and 0.36 for the strongest mean rule. It saves 21.92, 45.28, and 72.84 rounds and gains 21.22, 49.00, and 81.41 reward units; paired success intervals lie above zero. Total costs are higher (28.05, 121.14, and 303.62 versus 24.10, 106.71, and 271.00), so the advantage is higher success, speed, reward, and fairness rather than minimum cost. Reading the metrics jointly is important in this experiment. The rule baselines use the currently visible geometry and therefore react similarly when two members have comparable positions, even if their hidden responses differ. The proposed coordinator can revise the next suggestion after observing that one member accepted a previous recommendation while another member with a similar center did not. This adaptive probing increases early adjustment in some scenarios, which helps explain the additional total cost, but it also prevents repeated ineffective recommendations and sharply reduces termination time as $N$ grows. The result is thus a tradeoff rather than an unqualified dominance claim: the model spends more aggregate adjustment effort to preserve near-universal completion and distribute that effort more evenly.

\FloatBarrier
\begin{figure}[!t]
\centering
\subfloat[Success]{\includegraphics[width=0.76\columnwidth]{fig04a_success_rate.png}}\par\vspace{-0.6ex}
\subfloat[Rounds]{\includegraphics[width=0.76\columnwidth]{fig04b_convergence_steps.png}}\par\vspace{-0.6ex}
\subfloat[Reward]{\includegraphics[width=0.76\columnwidth]{fig04c_episode_reward.png}}\par\vspace{-0.6ex}
\subfloat[Cost]{\includegraphics[width=0.76\columnwidth]{fig04d_adjustment_cost.png}}
\caption{Hidden heterogeneity: (a) success; (b) rounds; (c) reward; and (d) cost.}
\label{fig:heterogeneity}
\end{figure}

\subsection{Context Ablation}

The zero-history condition masks only the ten-round encoder input while keeping the frozen actor, attention module, environment, and paired seeds unchanged. Real history raises success by 14, 22, and 44 percentage points at $N=10$, 40, and 100, with 95\% intervals $[6,24]$, $[12,34]$, and $[30,58]$. AUC gains are 0.0747, 0.0704, and 0.1338. Mean rounds fall from 35.74 to 17.40, 53.06 to 17.90, and 74.74 to 20.32; total cost falls by 7.45, 56.33, and 171.13. The paired design therefore attributes the difference to usable historical information, while a separately retrained context-free network remains a future structural ablation. The ablation also clarifies what the context encoder contributes. It does not add new members, reveal the hidden response parameters, or alter the environment; it changes only whether the frozen coordinator can consult the recorded interaction sequence. The larger success loss at $N=100$ indicates that a small per-member ambiguity can accumulate into a group-level coordination problem when many responses must be managed simultaneously. The simultaneous reductions in rounds and cost with real history show that the encoder is not merely making the actor more aggressive. Instead, previously observed acceptance and movement allow later proposals to be targeted with fewer unproductive trials. Because both conditions share initial scenarios and random seeds, these differences are evaluated within scenario rather than inferred from unrelated sample averages.

\begin{figure}[!t]
\centering
\subfloat[Success]{\includegraphics[width=0.76\columnwidth]{fig05a_success_rate.png}}\par\vspace{-0.6ex}
\subfloat[Rounds]{\includegraphics[width=0.76\columnwidth]{fig05b_convergence_steps.png}}\par\vspace{-0.6ex}
\subfloat[Cost]{\includegraphics[width=0.76\columnwidth]{fig05c_adjustment_cost.png}}\par\vspace{-0.6ex}
\subfloat[Reward]{\includegraphics[width=0.76\columnwidth]{fig05d_episode_reward.png}}
\caption{Context-history ablation: (a) success; (b) rounds; (c) cost; and (d) reward.}
\label{fig:context}
\end{figure}

\subsection{Interval Preservation and Scope}

Because Eq.~\eqref{eq:interval_update} translates both endpoints equally, Hausdorff distortion equals center displacement when $h_{i,T}=h_{i,0}$, as shown in Eq.~\eqref{eq:hausdorff}. Final center consensus is 0.9652, 0.9542, and 0.9482; mean half-width change is zero, and the maximum Hausdorff-center discrepancy is $1.11\times10^{-16}$. Fig.~\ref{fig:preservation} confirms that consensus is not obtained by narrowing members' uncertainty ranges. The slight decrease in final center consensus from $N=10$ to $N=100$ is therefore not evidence that uncertainty has been discarded. It reflects the more demanding task of aligning a larger collection of centers while every submitted half-width remains intact. This distinction matters because a center-only evaluation could report improvement after intervals were artificially contracted. Reporting zero mean half-width change together with the Hausdorff-center identity rules out that explanation and links the numerical success measure to the intended interval semantics.

\begin{equation}
\begin{aligned}
d_H(I_{i,T},I_{i,0})
&=\max\bigl(|l_{i,T}-l_{i,0}|,|u_{i,T}-u_{i,0}|\bigr)\\
&=|c_{i,T}-c_{i,0}|,\\
h_{i,T}&=h_{i,0}.
\end{aligned}
\label{eq:hausdorff}
\end{equation}

\eqnote{For intervals, Hausdorff distance is the larger endpoint displacement. Equal endpoint translation makes it equal to center shift while preserving $h_{i,t}$, verifying that consensus was not created by shrinking uncertainty.}

\begin{figure}[!t]
\centering
\subfloat[Center consensus]{\includegraphics[width=0.88\columnwidth]{fig06a_center_consensus.png}}\par\vspace{-0.6ex}
\subfloat[Hausdorff consistency]{\includegraphics[width=0.88\columnwidth]{fig06b_hausdorff_consistency.png}}
\caption{Interval preservation: (a) final center consensus and (b) Hausdorff-center consistency.}
\label{fig:preservation}
\end{figure}

The evidence is limited to one training seed, 50 frozen scenarios per cell, and half-widths up to 0.08. The architecture is context-conditioned meta-RL rather than MAML. Table~\ref{tab:sensitivity} makes the tested optimization and reachability boundaries explicit, while the following process evidence shows how those aggregate outcomes arise. Multiple training seeds and a matched retrained context-free model remain necessary next validation steps \cite{ref36}.

\subsection{Process Evidence and Reproducibility}

Code, the seed-20260507 checkpoint, frozen configurations, scenario outputs, figure inputs, claim mapping, and checksums are archived in the project repository. Fig.~\ref{fig:process} presents ordinary, polarized, and personality-response cases generated from the released numerical records. Polarized groups reach threshold in 13, 16, and 22 rounds at $N=10$, 40, and 100, with final consensus 0.9009, 0.9058, and 0.9003. In the ten-agent response case, consensus rises from 0.3936 to 0.9039 in five rounds. Every displayed trajectory preserves its interval band.

\begin{figure}[!t]
\centering
\subfloat[Ordinary $N=10$ trajectory]{\includegraphics[width=0.82\columnwidth]{figA1a_ordinary_n10.png}}\par\vspace{-0.6ex}
\subfloat[Polarized $N=100$ trajectory]{\includegraphics[width=0.82\columnwidth]{figA1b_polarized_n100.png}}\par\vspace{-0.6ex}
\subfloat[Personality-response consensus]{\includegraphics[width=0.82\columnwidth]{figA1c_persona_consensus.png}}\par\vspace{-0.6ex}
\subfloat[Acceptance probabilities]{\includegraphics[width=0.82\columnwidth]{figA1d_persona_acceptance.png}}
\caption{Process evidence: (a) ordinary negotiation; (b) polarized negotiation; (c) personality-response consensus; and (d) round-by-round acceptance probabilities.}
\label{fig:process}
\end{figure}

The ordinary trajectory begins from a dispersed but unimodal group. As suggestions are accepted, interval centers approach the group location while the upper and lower endpoints move in parallel. This is the trajectory-level counterpart of the Hausdorff result in Fig.~\ref{fig:preservation}: initial-to-final distance is explained by center translation rather than half-width reduction. The polarized case is more demanding because two center clusters begin far apart and members respond differently. The centralized actor therefore issues member-specific suggestions from the member state, encoded history, and social-attention representation rather than imposing one common displacement. The clusters move at different rates while retaining their submitted interval widths.

The response panels make the same mechanism visible at the individual level. Similar suggestions produce different realized movements because acceptance probabilities vary by member and round. This explains why historical context can reduce repeated probing under hidden heterogeneity, but the panels are mechanism illustrations rather than substitutes for the paired success, round, reward, cost, and fairness statistics. They also clarify the $\kappa=7$ boundary in Table~\ref{tab:sensitivity}: centers can become tightly concentrated while endpoint consensus remains below threshold because half-width differences are deliberately preserved. Forcing success would require narrowing the submitted uncertainty, relaxing the target, or changing the interval dynamics, none of which is done here.

\subsection{Reproduction and Reporting Scope}

A reproduction should first verify the released checksums and load the recorded configuration with the seed-20260507 checkpoint. Training is rerun only when the training claim itself is examined. During every frozen-policy comparison, the encoder, actor, critics, target networks, and normalization constants remain unchanged. Each method receives the same scenario file, including initial intervals, safety bounds, hidden response types, stubbornness, and cost sensitivities. Tests at $N=40$ and $N=100$ reuse the $N=10$ checkpoint; further fitting would answer a different question from scale transfer.

Scenario-level success, termination round, nonlinear cost, cost Gini, boundary violations, and cumulative reward are retained before aggregation. Failed runs keep $H=120$ in the round summary. Paired bootstrap intervals are recomputed from within-scenario differences, and figures are regenerated from stored records rather than copied from manuscript tables. The claim map links the training curves, rule and reinforcement-learning comparisons, hidden-heterogeneity test, history ablation, interval-preservation check, and process trajectories to their source files.

The released evidence has a deliberately limited scope. The main policy is trained with one random seed, and each evaluation cell contains 50 scenarios. History masking tests whether the frozen policy uses context, but it is not equivalent to training a separately optimized context-free architecture. The response-agent case uses fixed numerical mechanisms and is neither participant data nor an external large-language-model experiment. Multi-seed training, a matched retrained ablation, broader half-width distributions, and real negotiation records are appropriate next validation steps; none is implied by the present results.

\FloatBarrier
\section{Conclusion}
\label{sec:conclusion}

Interval Meta-GDM couples endpoint-consistent consensus control with history-conditioned response inference and set-based action generation. The model does not dominate simple rules in static homogeneous tasks, but it retains their success while offering better composite reward and burden balance at larger scales. Its clearest advantage appears when hidden response differences accumulate: the frozen policy uses observed interaction history to sustain success, reduce repeated probing, and transfer from $N=10$ to $N=40$ and $N=100$ without architectural changes.


From an application perspective, the proposed system should be understood as decision support for a moderator rather than an autonomous replacement for participants. The coordinator recommends bounded movements after observing the group state, but individuals retain their hidden acceptance behavior and may realize only part of a suggestion. Preserving the submitted half-width makes this separation visible: the model can move an opinion interval toward agreement without rewriting the uncertainty range that the member declared acceptable. The recorded history further provides an audit trail linking each later suggestion to earlier proposals, realized movements, and rewards. These properties are useful when a practitioner must explain why two similarly positioned members received different recommendations. They do not, however, remove the need for domain-specific approval of the safety bounds, consensus threshold, and cost interpretation before deployment.

The interval extension is structural rather than cosmetic: actions preserve half-widths, boundary checks use both endpoints, and the reward uses endpoint consensus, cost, fairness, and violations. The main remaining issues are higher total cost under heterogeneity, single-seed training, and incomplete structural ablation. Future work should optimize this trade-off, repeat training across seeds, and test task-level adaptation and real negotiation data.

\FloatBarrier
\begin{thebibliography}{39}
\setlength{\itemsep}{-0.45ex}

\bibitem{ref1} L. A. Zadeh, ``Fuzzy sets,'' \emph{Information and Control}, vol. 8, no. 3, pp. 338--353, 1965.
\bibitem{ref6} M. Cao, J. Wu, F. Chiclana, R. Ure\~na, and E. Herrera-Viedma, ``A personalized consensus feedback mechanism based on maximum harmony degree,'' \emph{IEEE Transactions on Systems, Man, and Cybernetics: Systems}, vol. 51, no. 10, pp. 6134--6146, 2021.
\bibitem{ref25} X. Yao and Z. Xu, ``A survey of consensus in group decision making under the CWW environment,'' \emph{Applied Soft Computing}, vol. 144, Art. no. 110557, 2023.
\bibitem{tfsWu2022} T. Wu, X. Liu, J. Qin, and F. Herrera, ``Trust-consensus multiplex networks by combining trust social network analysis and consensus evolution methods in group decision-making,'' \emph{IEEE Transactions on Fuzzy Systems}, vol. 30, no. 11, pp. 4741--4753, 2022.
\bibitem{tfsLiu2023} P. Liu, Y. Li, and P. Wang, ``Social trust-driven consensus reaching model for multiattribute group decision making: Exploring social trust network completeness,'' \emph{IEEE Transactions on Fuzzy Systems}, vol. 31, no. 9, pp. 3040--3054, 2023.
\bibitem{tfsJi2024} F. Ji, J. Wu, F. Chiclana, Q. Sun, C. Liang, and E. Herrera-Viedma, ``Decayed trust propagation method in multiple overlapping communities for improving consensus under social network group decision making,'' \emph{IEEE Transactions on Fuzzy Systems}, vol. 32, no. 8, pp. 4401--4412, 2024.
\bibitem{ref16} L. Zhang, Z. Yang, and T. Li, ``Group decision making with incomplete interval-valued linguistic intuitionistic fuzzy preference relations,'' \emph{Information Sciences}, vol. 647, Art. no. 119451, 2023.
\bibitem{ref17} Z. Zhang and S.-M. Chen, ``Optimization-based group decision making using interval-valued intuitionistic fuzzy preference relations,'' \emph{Information Sciences}, vol. 561, pp. 352--370, 2021.
\bibitem{ref2} M. H. DeGroot, ``Reaching a consensus,'' \emph{Journal of the American Statistical Association}, vol. 69, no. 345, pp. 118--121, 1974.
\bibitem{ref21} C. Shang, R. Zhang, X. Zhu, and Y. Liu, ``An adaptive consensus method based on feedback mechanism and social interaction in social network group decision making,'' \emph{Information Sciences}, vol. 625, pp. 430--456, 2023.
\bibitem{ref27} L. Li, S. Jiao, Y. Shen, \emph{et al.}, ``A two-stage consensus model for large-scale group decision-making considering dynamic social networks,'' \emph{Information Fusion}, vol. 100, Art. no. 101972, 2023.
\bibitem{ref7} J. Du, S. Liu, and Y. Liu, ``A limited cost consensus approach with fairness concern and its application,'' \emph{European Journal of Operational Research}, vol. 298, no. 1, pp. 261--275, 2022.
\bibitem{ref8} G. Gong, K. Li, and Q. Zha, ``A maximum fairness consensus model with limited cost in group decision making,'' \emph{Computers \& Industrial Engineering}, vol. 175, Art. no. 108891, 2023.
\bibitem{ref20} J. Guo, Z. Wang, and Z. Zhang, ``Minimum cost consensus model with variable cost for multi-criteria large-scale group decision-making considering the tolerance of decision-makers,'' \emph{Computers \& Industrial Engineering}, vol. 196, Art. no. 110456, 2024.
\bibitem{ref12} X. Liu, Y. Shan, Y. Xu, and X. Wang, ``Minimum cost consensus modeling considering opinion-based variable cost and trust-based tolerance constraints,'' \emph{IEEE Transactions on Computational Social Systems}, vol. 12, no. 5, pp. 3023--3036, 2025.
\bibitem{ref26} T. Gai, J. Wu, M. Cao, F. Ji, Q. Sun, and M. Zhou, ``Trust chain driven bidirectional feedback mechanism in social network group decision making and its application in Metaverse virtual community,'' \emph{Expert Systems with Applications}, vol. 228, Art. no. 120369, 2023.
\bibitem{tfsLiang2022} H. Liang, G. Kou, Y. Dong, F. Chiclana, and E. Herrera-Viedma, ``Consensus reaching with minimum cost of informed individuals and time constraints in large-scale group decision-making,'' \emph{IEEE Transactions on Fuzzy Systems}, vol. 30, no. 11, pp. 4991--5004, 2022.
\bibitem{tfsZhang2022} H. Zhang, F. Wang, Y. Dong, F. Chiclana, and E. Herrera-Viedma, ``Social trust driven consensus reaching model with a minimum adjustment feedback mechanism considering assessments-modifications willingness,'' \emph{IEEE Transactions on Fuzzy Systems}, vol. 30, no. 6, pp. 2019--2031, 2022.
\bibitem{tfsDai2024} X. Dai, H. Li, L. Zhou, Q. Wu, W. Ding, and M. Deveci, ``A Bayesian framework for modelling the trust relationships to group decision making problems,'' \emph{IEEE Transactions on Fuzzy Systems}, vol. 32, no. 12, pp. 6594--6606, 2024.
\bibitem{ref28} F. Jin, Y. Guan, J. Liu, and L. Zhou, ``Bayesian inference and minimum consensus adjustment process for multi-attribute large-scale group decision-making in social networks,'' \emph{Information Fusion}, vol. 126, Art. no. 103537, 2026.
\bibitem{ref29} F. Jin, Z. Li, and L. Zhou, ``Managing group decision-making with bounded rationality: A robust behavioral framework for consensus building and psychological preference analysis,'' \emph{Computers \& Industrial Engineering}, vol. 217, Art. no. 112036, 2026.
\bibitem{ref30} M. Tang, M. Li, and L. Xie, ``Harnessing the hybrid intelligence of crowd and artificial intelligence in group decision making for uncertain disaster response,'' \emph{Risk Analysis}, vol. 46, no. 5, Art. no. e70254, 2026.
\bibitem{ref14} Y. Shen, X. Ma, Y. Bao, G. Kou, and J. Zhan, ``A consensus method based on reinforcement learning for group decision-making,'' \emph{European Journal of Operational Research}, vol. 330, no. 3, pp. 941--955, 2026.
\bibitem{ref15} Y. Tu, W. Wang, Z. Ma, Z. Li, and B. Lev, ``Multi-agent reinforcement learning-based consensus building for large-scale group decision making,'' \emph{Information Fusion}, vol. 126, Art. no. 103629, 2026.
\bibitem{ref18} Y. Xu and H. Xu, ``A maximum satisfaction consensus-based large-scale group decision-making in social network considering limited compromise behavior,'' \emph{Information Sciences}, vol. 670, Art. no. 120606, 2024.
\bibitem{ref23} W. Wang, Z. Liu, X. Feng, and P. Liu, ``A consensus model for managing short- and long-term non-cooperative behaviors in social network group decision-making: The perspective of historical and current performance,'' \emph{Expert Systems with Applications}, vol. 266, Art. no. 126112, 2025.
\bibitem{ref24} H. Zhang, J. Wang, and W. Zhu, ``Modeling consistency and consensus in social network group decision making: The role of limited dual tolerance and compromise behaviors,'' \emph{Applied Soft Computing}, vol. 166, Art. no. 112130, 2024.
\bibitem{ref9} R. Zhang, J. Huang, Y. Xu, and E. Herrera-Viedma, ``Consensus models with aggregation operators for minimum quadratic cost in group decision making,'' \emph{Applied Intelligence}, vol. 53, no. 2, pp. 1370--1390, 2023.
\bibitem{ref10} Y. Shen, X. Ma, Z. Xu, M. Deveci, and J. Zhan, ``A minimum cost and maximum fairness-driven multi-objective optimization consensus model for large-scale group decision-making,'' \emph{Fuzzy Sets and Systems}, vol. 500, Art. no. 109198, 2025.
\bibitem{ref11} Y. Liang, J. Qin, and W. Pedrycz, ``Minimum cost consensus model considering dual behavior preference,'' \emph{Computers \& Operations Research}, vol. 176, Art. no. 106961, 2025.
\bibitem{ref13} B. Zhang, W. Huang, Y. Dong, and W. Pedrycz, ``Utility-driven minimum cost consensus model with preference learning and fairness concerns,'' \emph{IEEE Transactions on Systems, Man, and Cybernetics: Systems}, vol. 55, no. 7, pp. 4641--4655, 2025.
\bibitem{ref3} A. Vaswani, N. Shazeer, N. Parmar, \emph{et al.}, ``Attention is all you need,'' in \emph{Advances in Neural Information Processing Systems}, vol. 30, pp. 5998--6008, 2017.
\bibitem{ref31} P. Verbeke and T. Verguts, ``Reinforcement learning and meta-decision-making,'' \emph{Current Opinion in Behavioral Sciences}, vol. 57, Art. no. 101374, 2024.
\bibitem{ref34} S. Arora and P. Doshi, ``A survey of inverse reinforcement learning: Challenges, methods and progress,'' \emph{Artificial Intelligence}, vol. 297, Art. no. 103500, 2021.
\bibitem{ref4} J. Lee, Y. Lee, J. Kim, \emph{et al.}, ``Set Transformer: A framework for attention-based permutation-invariant neural networks,'' \emph{Proceedings of Machine Learning Research}, vol. 97, pp. 3744--3753, 2019.
\bibitem{ref19} Y. Xu, Y. Ju, Z. Gong, \emph{et al.}, ``Improving consensus in social network group decision-making: Emphasizing overlapping subgroups and interactive behaviors,'' \emph{Information Sciences}, vol. 679, Art. no. 121065, 2024.
\bibitem{ref5} T. Haarnoja, A. Zhou, P. Abbeel, and S. Levine, ``Soft actor--critic: Off-policy maximum entropy deep reinforcement learning with a stochastic actor,'' \emph{Proceedings of Machine Learning Research}, vol. 80, pp. 1861--1870, 2018.
\bibitem{ref22} Y. Xiao, X. Ma, J. C. R. Alcantud, and J. Zhan, ``A group consensus method based on social network and three-way decision under multi-scale information systems,'' \emph{Applied Soft Computing}, vol. 162, Art. no. 111824, 2024.
\bibitem{ref36} T. Rupprecht and Y. Wang, ``A survey for deep reinforcement learning in Markovian cyber--physical systems: Common problems and solutions,'' \emph{Neural Networks}, vol. 153, pp. 13--36, 2022.

\end{thebibliography}

\end{document}


